%%TOPSIS
%%Technique for Order Preference by Similarity to an Ideal Solution
%%理想解法 优劣解距离法


\subsection{TOPSIS}
Here is problem of evaluating the water quality of 20 rivers.

\subsubsection{Data set}
Four indexes have been set and measured in the following table (see in Table \ref{tbl:TOPSIS-water quality data set}).
%%Data set table
%%三线表
\begin{table}[H]
    \renewcommand\arraystretch{2}   %%控制行间距
    \centering
    \begin{tabular}{c r@{.}l r@{.}l r@{.}l c}    %%各列对齐格式设置
        \toprule  %%表格横线加粗
        \specialrule{0em}{2pt}{2pt}    %%添加透明线控制行高
        \textbf{River}& \multicolumn{2}{l}{\textbf{Oxygen content}} & \multicolumn{2}{l}{\textbf{pH}} & \multicolumn{2}{l}{\textbf{Bacteria number}} & \textbf{Plant nutrients} \\
        \specialrule{0em}{2pt}{2pt}    %%添加透明线控制行高
        \midrule
        \specialrule{0em}{1pt}{1pt}
        1 & 4&69 & 6&59 & 51 & 11&94 \\
        2 & 2&03 & 7&86 & 19 & 6&46 \\
        3 & 9&11 & 6&31 & 46 & 8&91 \\
        4 & 8&61 & 7&05 & 46 & 26&43 \\
        5 & 7&13 & 6&50 & 50 & 23&57 \\
        6 & 2&39 & 6&77 & 38 & 24&62 \\
        7 & 7&69 & 6&79 & 38 & 6&01 \\
        8 & 9&30 & 6&81 & 27 & 31&57 \\
        9 & 5&45 & 7&62 & 5 & 18&46 \\
        10 & 6&19 & 7&27 & 17 & 7&51 \\
        11 & 7&93 & 7&53 & 9 & 6&52 \\
        12 & 4&40 & 7&28 & 17 & 25&30 \\
        13 & 7&46 & 8&24 & 23 & 14&42 \\
        14 & 2&01 & 5&55 & 47 & 26&31 \\
        15 & 2&04 & 6&40 & 23 & 17&91 \\
        16 & 7&73 & 6&14 & 52 & 15&72 \\
        17 & 6&35 & 7&58 & 25 & 29&46 \\ 
        18 & 8&29 & 8&41 & 39 & 12&02 \\
        19 & 3&54 & 7&27 & 54 & 3&16 \\
        20 & 7&44 & 6&26 & 8 & 28&41 \\
        \specialrule{0em}{1pt}{1pt}
        \bottomrule
    \end{tabular}
    \caption{TOPSIS-water quality data set.}
    \label{tbl:TOPSIS-water quality data set}
\end{table}


%%翻译不确定 极大型指标 极小型指标 中间型指标 区间型指标
\subsubsection{Data conversion}
In this part, the four indexes should be processed to 'very large index'.

(a)Oxygen content
High oxygen content indicates better water quality, so this index is a 'very large index', which doesn't need any processing.

(b)pH
Since the pH of water should be near 7, we have a 'middle index', and the best pH is 7.

To convert a 'middle index' to 'very large index', the following equation is applied here.

\begin{equation*}
    M = \max {|x_i-x_{best}|} , x_i' = 1 - \frac{|x_i-x_{best}|}{M}
\end{equation*}

(c)Bacteria number
The better the water quality is, the fewer the bacteria number should be. So the bacteria number is a 'very little index', which can be converted through the following way.

\begin{equation*}
    x_i' = max - x_i
\end{equation*}

(d)Plant nutrient
The index of the plant nutrient have an appropriate interval between 10 to 20, which should be converted to 'very large index' by using the following way.

In an 'interval index', $[a,b]$ is the stable interval and $[a^*,b^*]$ is the tolerant interval (the min and the max of data).

\begin{equation*}
    x' = \begin{cases}
    1 - \frac{a-x}{a-a*}, & x<a ; \\
    1 &, a \leq x \leq b ; \\
    1 - \frac{x-b}{b*-b}, & x>b .
    \end{cases}
\end{equation*}

\subsubsection{Construct the normalized initial matrix}
Then the normalized initial matrix can be constructed.

\subsubsection{Construct the index weight by AHP}

\subsubsection{Construct the evaluation system}